
A hallucination detector calibrated on TriviaQA may exhibit substantially different performance when applied to other benchmarks---degrading by approximately 22\% on PopQA versus approximately 3\% on SQuAD in the experiments reported here. This variability reveals a gap in understanding when uncertainty-based detection methods transfer across benchmarks.

Large language models produce confident but incorrect responses at rates that complicate deployment in high-stakes applications. Semantic entropy, which measures uncertainty by clustering semantically equivalent generations and computing entropy over these clusters, has emerged as a detection method achieving approximately 0.85 AUROC on factual QA benchmarks \citep{kuhn2023semantic}. However, calibration thresholds tuned on one benchmark may not generalize to others, and prior work has not established criteria for predicting which benchmark pairs exhibit compatible uncertainty distributions.

This gap has practical implications. A medical QA system calibrated on one benchmark may encounter queries with different error-generation characteristics. Without transfer criteria, each new domain requires recalibration---a cost that could potentially be reduced if practitioners could predict which benchmark pairs share compatible distributions.

Prior work has evaluated benchmarks independently, implicitly treating ''factual QA'' as a homogeneous category. The present investigation tests whether this assumption holds. The central hypothesis is that benchmarks cluster into families based on uncertainty distribution similarity and that transfer success depends on cluster membership.

This paper presents a systematic investigation of cross-benchmark transfer for uncertainty-based hallucination detectors. The approach uses Jensen-Shannon divergence to quantify distribution similarity between benchmark pairs, hierarchical clustering to discover benchmark families, and a calibration-transfer protocol to measure degradation when applying source-trained thresholds to target benchmarks.

The experiments across six benchmarks reveal:

\begin{enumerate}
\item \textbf{Benchmark taxonomy via uncertainty distributions.} QA benchmarks cluster into two families---Factual Recall (TriviaQA, NQ, SQuAD) and Entity/Claim Verification (PopQA, HaluEval, FEVER)---with silhouette score 0.8245.
\end{enumerate}

\begin{enumerate}
\item \textbf{Quantified transfer boundaries.} Within-cluster threshold transfer shows mean AUROC degradation of 0.032, while cross-cluster transfer shows degradation of 0.223---a ratio of approximately $7\times$.
\end{enumerate}

\begin{enumerate}
\item \textbf{Cluster membership as transfer criterion.} JS-divergence clustering provides a candidate criterion for predicting transfer success prior to deployment.
\end{enumerate}

